\section*{Abstract}
Optical character recognition engines always return a transcription, so every pipeline must decide, word by word,
whether to trust it. Existing verifiers judge the claim in isolation, by
engine confidence, by agreement with a second recogniser, or by asking a large vision-language model. We instead ask
whether the claim is the best explanation of the image among the strings it could have been confused with. LOCUS builds around each claim a small graph of plausible single-character edits and independent
readings, scores every candidate string by how much better it explains the real image than a clean rendering of
itself, cancelling the recogniser's language prior, and tests whether the claim is locally optimal with a
relation-aware graph neural network. The network is pretrained without human labels from engine agreement and
synthetic corruptions and adapts from a few labelled words. We introduce a benchmark of real recognition errors from
two engines on five public datasets covering forms, receipts, handwriting and scene text. With 25 labelled words per
class, LOCUS ranks correct above incorrect transcriptions better than engine confidence, learned confidence, engine
voting, Florence-2 re-recognition and a seven-billion-parameter vision-language judge on every dataset, and attains
the best average macro-F1. It is significantly better than re-recognition with a strong recogniser
where that recogniser is fallible, comparable where it is in-domain, and behind it on its own scene-text benchmark.
Significance, recogniser, rendering and message-passing analyses and a manual audit show that the gain comes from explicit, prior-cancelled counterfactual evidence and label-free pretraining.
\vspace{5pt}
\keywords{OCR verification; counterfactual reasoning; graph neural networks; few-shot learning; render normalisation; document analysis}
